\begin{thebibliography}{9}
\raggedright
\bibitem{aht}
Ian Agol, Joel Hass, and William Thurston.
\emph{The computational complexity of knot genus and spanning area}.
\href{https://arxiv.org/abs/math/0205057}{arXiv:math/0205057}.
See Theorems 12 and 16, Corollary 17, and the normal-coordinate
orientability test in Section 5.

\bibitem{lackenby-talk}
Marc Lackenby.
\emph{Unknot recognition in quasi-polynomial time}.
Oxford lecture slides, 2021, especially slides 12--13, 72--73, and 108.
\url{https://people.maths.ox.ac.uk/lackenby/quasipolynomial-talk-oxford.pdf}.

\bibitem{lackenby2026}
Marc Lackenby.
\emph{Incompressible surfaces, hierarchies and unknot recognition}.
\href{https://arxiv.org/abs/2607.23350}{arXiv:2607.23350v1}, July 2026.
See Sections 9--10 for the quantitative scope of the described algorithms.

\bibitem{lackenby-curves}
Marc Lackenby.
\emph{Some fast algorithms for curves in surfaces}.
\emph{Discrete \& Computational Geometry} 76 (2026), 539--588.
\href{https://doi.org/10.1007/s00454-026-00845-7}{doi:10.1007/s00454-026-00845-7}.
Author manuscript:
\url{https://people.maths.ox.ac.uk/lackenby/AlgorithmsSurfaces13jan26.pdf}.

\bibitem{regina}
The Regina development team.
\emph{Regina: software for low-dimensional topology}.
NormalSurface API documentation, especially \code{components},
\code{countBoundaries}, and \code{isOrientable}.
\url{https://regina-normal.github.io/engine-docs/classregina_1_1NormalSurface.html}.
The external oracle in this package uses installed Regina 7.4.

\bibitem{proveit-baseline}
Vladimir Reshetnikov, ProveIt contributors, and prior research continuations.
\emph{ProveIt: Topology/UnknotRecognition}.
Pinned source commit \baseline.
\url{https://github.com/VladimirReshetnikov/ProveIt/tree/eb368edf975695e3e16a8774dcb7846bda0c13a0/Topology/UnknotRecognition}.
Relevant maintained chapters: \path{normal_multiplicity.tex},
\path{normal_boundary.tex}, \path{weighted_components.tex}, and
\path{surface_covers.tex}; exact runtime and research snapshot hashes are
included in the accompanying package.
\end{thebibliography}
